\begin{pastexam}[label={px:ncu-111-dsa-14}]{中央資工所 111 學年度 問答題第 2 題}
（問答）

(13\%) You are going on a long hiking trip. You start on the road at kilometer post 0. Along the way
there are $n$ hotels, at kilometer posts $a_1 < a_2 < \cdots < a_n$, where each $a_i$ is measured
from the starting point. The only places you are allowed to stop are at these hotels, but you can
choose which of the hotels you stop at. You must stop at the final hotel (at distance $a_n$), which
is your destination. You'd ideally like to travel 20 kilometers a day, but this may not be
possible. If you travel $x$ kilometers during a day, the penalty for that day is $(20-x)^2$. You
want to plan your trip so as to minimize the total penalty. Give an algorithm that determines the
optimal sequence of hotels at which to stop. (Note: This problem can be solved by a dynamic
programming algorithm. If you use a dynamic programming algorithm to solve this problem, please
write down the recursive formula of the algorithm, describe the meaning of the notations you used
in the formula, give the initial value settings, and analyze the time complexity of the algorithm.)
\end{pastexam}

\begin{pastsoln}
照資結書第~18.1.5~節的步驟：定義子問題、寫遞迴式、給初值、決定填表順序、記錄選擇以便還原答案。

\textbf{記號。} 令 $a_0 = 0$（起點）。 $P[j]$（$0 \le j \le n$）表示「從起點出發、這一天結束時停在旅館
$j$」的所有走法裡，到目前為止的總懲罰的最小值。 $\mathit{prev}[j]$ 記錄達到 $P[j]$ 時，前一晚停在哪裡。

\textbf{遞迴式。} 最後一天是從某個 $i < j$ 走到 $j$，走 $a_j - a_i$ 公里，懲罰 $(20 - (a_j - a_i))^2$；
在那之前是一個停在 $i$ 的走法，它本身也必須是最佳的（否則換成停在 $i$ 的最佳走法，總懲罰會更小）。 所以
\[
P[j] = \min_{0 \le i < j} \Bigl( P[i] + \bigl(20 - (a_j - a_i)\bigr)^2 \Bigr), \qquad 1 \le j \le n,
\]
$\mathit{prev}[j]$ 是取到最小值的 $i$。

\textbf{初值。} $P[0] = 0$：還沒出發，沒有懲罰。 答案是 $P[n]$（一定要停在最後一家）。

\begin{pseudo}{旅館停靠的最小懲罰}
\Function{Hotel-Stops}{$a, n$} \Comment{$a[1..n]$ 遞增}
    \State $a[0] \gets 0$;\ $P[0] \gets 0$
    \For{$j \gets 1$ \To $n$}
        \State $P[j] \gets \infty$
        \For{$i \gets 0$ \To $j - 1$}
            \State $c \gets P[i] + (20 - (a[j] - a[i]))^2$
            \If{$c < P[j]$}
                \State $P[j] \gets c$;\ $\mathit{prev}[j] \gets i$
            \EndIf
        \EndFor
    \EndFor
    \State $S \gets$ 空串列;\ $j \gets n$
    \While{$j > 0$}
        \State 把 $j$ 加到 $S$ 的最前面;\ $j \gets \mathit{prev}[j]$
    \EndWhile
    \State \Return $P[n]$ 和停靠的旅館序列 $S$
\EndFunction
\end{pseudo}

$P[j]$ 只用到 $P[0..j-1]$，所以 $j$ 由小到大填。 還原時從 $n$ 沿 $\mathit{prev}$ 往回走到 $0$，倒過來就是
要停的旅館。

\textbf{例子。} 旅館在 $12, 21, 33, 40$ 公里：$P[1] = 8^2 = 64$；$P[2] = \min(0 + 1^2,\; 64 + 11^2) = 1$；
$P[3] = \min(0 + 13^2,\; 64 + 1^2,\; 1 + 8^2) = 65$；$P[4] = \min(0 + 20^2,\; 64 + 8^2,\; 1 + 1^2,\;
65 + 13^2) = 2$。 最佳是停 $21$、$40$ 兩家，總懲罰 $2$。

\textbf{時間。} 外層 $n$ 輪，第 $j$ 輪的內層做 $j$ 次 $O(1)$ 的計算，共
$\sum_{j=1}^{n} j = n(n+1)/2 = \Theta(n^2)$；還原 $O(n)$。 總時間 $\Theta(n^2)$，空間 $\Theta(n)$。
\end{pastsoln}
